\documentclass[handout, font=12pt]{ximera}

\title{Fractions: Introduction}   % Use xmlatex all -s to compile when SERVE refuses to work.
\author{Kelly Stady}
\license{CC: 0}         % replace with an appropriate license, or set it in xmPreamble

\begin{document}   

\begin{abstract}

\end{abstract}

\maketitle

\label{xim:FractionsIntro}

\section*{Introduction to Fractions}

A \textbf{fraction} is an \emph{equal} part of a whole. The square below has been divided into 4 equal parts.  The fraction 
that represents the shaded part of the square is $\textstyle \frac{1}{4},$ since 1 of the 4 parts is shaded.

\begin{center}
\begin{tikzpicture}
  % 1. Shade the top-right quadrant (0,0 to 1.3,1.3)
  %\fill[glaucous!60] (0,0) -- (1.3,0) -- (1.3,1.3) -- (0,1.3) -- cycle;

  % 1a. Shade the top-left quadrant
  \fill[glaucous!60] (0,0) rectangle (-1.3,1.3);

  % 2. Draw the main square outline (centered at 0,0 like your circle was)
  \draw[thick] (-1.3,-1.3) rectangle (1.3,1.3);

  % 3. Draw the division lines (horizontal and vertical through the center)
  \draw[thick] (-1.3,0) -- (1.3,0);
  \draw[thick] (0,-1.3) -- (0,1.3);

  % Label
    \node at (0,-2) {\Large $\frac{1}{4}$};
\end{tikzpicture}

%Alt Text
\xmalt{A square divided into four equal quadrants by a horizontal and a vertical line through its center. 
The top-left quadrant is shaded blue, representing the fraction one-fourth.}

\end{center}

The \textbf{denominator} of a fraction tells you how many equal-sized pieces a whole has been divided into.  The \textbf{numerator} 
of a fraction is how many of those pieces you are counting or considering.


\textbf{Why do we need to divide a whole into equal parts}?

Because fractions only make sense if the whole is divided into equal parts.  For example, the INCORRECT square on the left has been divided 
into 3 parts and 1 of them is shaded.  However, we can't say that $\textstyle \frac{1}{3}$ of the square is shaded because the parts are not 
the same size.  To correct this, we redraw the lines so that every part is the exact same size, as shown in the CORRECT square on the right.


\begin{center}
\begin{tikzpicture}[thick]
  % --- INCORRECT FIGURE (Left) ---
  \begin{scope}
    \fill[glaucous!60] (-1.3,-1.3) rectangle (0.3,1.3);
    \draw[ultra thick] (-1.3,-1.3) rectangle (1.3,1.3);
    \draw (0.3,-1.3) -- (0.3,1.3);
    \draw (0.3,0) -- (1.3,0);
    \node at (0,-2) {\Large $\frac{1}{3}$ ?};

    % INCORRECT Header and X (Stacked)
    \node[accessibleRed] at (0, 2.3) {\textbf{INCORRECT}};
    % Shifted down slightly to 1.6/2.0
    \draw[accessibleRed, ultra thick] (-0.2, 1.6) -- (0.2, 2.0);
    \draw[accessibleRed, ultra thick] (0.2, 1.6) -- (-0.2, 2.0);
  \end{scope}

  % --- CORRECT FIGURE (Right) ---
  \begin{scope}[xshift=4.5cm]
    \fill[glaucous!60] (-1.3,-1.3) rectangle (-0.433,1.3);
    \draw[ultra thick] (-1.3,-1.3) rectangle (1.3,1.3);
    \draw (-0.433,-1.3) -- (-0.433,1.3);
    \draw (0.433,-1.3) -- (0.433,1.3);
    \node at (0,-2) {\Large $\frac{1}{3}$};

    % CORRECT Header and Checkmark (Stacked)
    \node[accessibleGreen] at (0, 2.3) {\textbf{CORRECT}};
    \draw[accessibleGreen, ultra thick] (-0.2, 1.8) -- (0, 1.6) -- (0.3, 2.0);   
  \end{scope}
\end{tikzpicture}

% alt tex
 \xmalt{Two squares side-by-side demonstrating the importance of equal parts in fractions. The left square is 
 labeled INCORRECT with a red X; it is divided into three unequal parts with one large section shaded, showing 
 that counting unequal pieces as one-third is incorrect. The right square is labeled CORRECT with a green checkmark; 
 it is divided into three identical vertical rectangles with one shaded, correctly representing the fraction one-third.}

\end{center}


\begin{problem}  The circle below has been divided into 10 equal parts.

\begin{center}
\begin{tikzpicture}
\draw[thick] (0,0) circle [radius=1.3cm];
\foreach \i in {36,72,...,360}
  \draw[thick] (0,0)--(\i:1.3cm);    % The colon signifies that polar coordinates are being used. 
                                     % (60:.75cm) means the point that is at an angle of 60° and a 
                                     % distance of 0.75cm from the origin. Length of radial lines - should match radius above.
 \draw[thick, fill=glaucous!60] (0,0) --  (36:1.3) arc(36:0:1.3);   % To Draw Arc: \draw (x,y) arc (start:stop:radius);
 \draw[thick, fill=glaucous!60] (0,0) --  (72:1.3) arc(72:36:1.3);  % Colon indicates polar coordintes (angle:length)
 \draw[thick, fill=glaucous!60] (0,0) --  (108:1.3) arc(108:72:1.3);
\end{tikzpicture}
%Alt Text
\xmalt{A circle representing a whole, divided into 10 equal pie-shaped wedges. Three adjacent wedges in the top-right section are
 shaded blue.}
\end{center}

\normalsize % In the html this font looks larger than the problem statement above, so trying to fix. 
What fraction of the circle is shaded?

\[ \textrm{Fraction Shaded } = \frac{\answer{3}}{\answer{10}}  \]

\end{problem}


\begin{problem}  A local wildlife sanctuary covers 12 acres of protected land.  Recent data shows that 5 acres have been successfully restored with native wild flowers to support declining bee populations.


\begin{center}
\begin{tikzpicture}[thick]
  % Define a simple flower shape 
  \newcommand{\flower}[2]{
    \foreach \a in {0,72,...,288}
      \draw[very thin, fill=white] (#1,#2) +(\a:0.24) circle (0.12);
    %\fill [orange!80!white] (#1,#2) circle (0.09);  %Use fill if you don't want the black outline
    % Use \draw[fill=...] to get BOTH the black outline and the orange inside
    \draw[very thin, fill=orange!80!white] (#1,#2) circle (0.09);
  }

  % Shade the 5 restored acres (Top row and first box of middle row)
  \fill[glaucous!60] (0, 2.6) rectangle (5.2, 3.9);
  \fill[glaucous!60] (0, 1.3) rectangle (1.3, 2.6);

  % Draw the flowers in the shaded boxes
  % Top row flowers
  \flower{0.65}{3.25}
  \flower{1.95}{3.25}
  \flower{3.25}{3.25}
  \flower{4.55}{3.25}
  
  % Middle row flower (first box)
  \flower{0.65}{1.95}

  % Draw the grid after to ensure lines are visible on top
  \draw (0,0) rectangle (5.2, 3.9);
  \foreach \x in {1.3, 2.6, 3.9} \draw (\x, 0) -- (\x, 3.9);
  \draw (0, 1.3) -- (5.2, 1.3);
  \draw (0, 2.6) -- (5.2, 2.6);

  \node at (2.6,-0.6) {\textbf{Each box represents 1 acre of the sanctuary}.};
\end{tikzpicture}

% Alt Text
\xmalt{A grid of 12 equal squares arranged in three rows of four. The four squares in the top row and the first square in
 the middle row are shaded blue with white flowers in the center.}
\end{center}

\begin{hint}
The denominator in each part is the number of parts the sanctuary has been divided into.  The numerator in the first part is
the number of squares with flowers. The numerator in the second part is the number of squares without flowers.
\end{hint}

What fraction of the sanctuary's total area has been restored with wild flowers?

\[ \frac{\answer{5}}{\answer{12}} \] 


What fraction of the sanctuary's total area has \textbf{not} been restored with wild flowers?

\[ \frac{\answer{7}}{\answer{12}} \] 

\end{problem}


\subsection*{Equivalent Fractions}

The square on the left has been divided into 4 equal parts, with 2 parts shaded.  The square on the right has been divided into 2 equal parts, with 
1 part shaded.  Since both squares are the same size and have an equal amount of shading, we can conclude that $\textstyle \frac{2}{4}$ is equivalent to
$\textstyle \frac{1}{2}.$


\begin{center}
\begin{tikzpicture}[thick]
  % --- LEFT SQUARE (2/4) ---
  % Shade the left half (both left quadrants)
  \fill[glaucous!60] (-1.3,-1.3) rectangle (0,1.3);

  % Division lines
  \draw (-1.3,0) -- (1.3,0); % Horizontal
  \draw (0,-1.3) -- (0,1.3); % Vertical

  % Outer Border
  \draw (-1.3,-1.3) rectangle (1.3,1.3);

  % Label
  \node at (0,-2) {\Large $\frac{2}{4}$};

  % --- EQUAL SIGN ---
  \node at (2.2, 0) {\Huge $=$};

  % --- RIGHT SQUARE (1/2) ---
  \begin{scope}[xshift=4.4cm]
    % Shade the left half
    \fill[glaucous!60] (-1.3,-1.3) rectangle (0,1.3);

    % Vertical division line
    \draw (0,-1.3) -- (0,1.3);

    % Outer Border
    \draw (-1.3,-1.3) rectangle (1.3,1.3);

    % Label
    \node at (0,-2) {\Large $\frac{1}{2}$};
  \end{scope}

\end{tikzpicture}

% Alt Text
\xmalt{Two squares side-by-side, as described above, separated by an equal sign.  The square on the left represents two-fourths and the square
on the right represents one-half.}

\end{center}


\textbf{Equivalent fractions} represent the same part of a whole, but have different numerators and denominators.  Since a whole can be divided
into any number of equal parts, each fraction has infinitely many equivalent fractions.


\begin{problem} Consider the circle below.

\begin{center}
\begin{tikzpicture}[thick]
  % Original 2/6 Shading
  \fill[glaucous!60] (0,0) -- (90:1.3) arc (90:210:1.3) -- cycle;
  \draw (0,0) circle (1.3);
  \foreach \a in {30, 90, 150, 210, 270, 330}
    \draw (0,0) -- (\a:1.3);
\end{tikzpicture}
% alt tex
 \xmalt{A circle divided into six equal wedge-shaped sections. Two adjacent sections on the upper-left are shaded blue.}
\end{center}


What fraction of the circle is shaded?

\[ \textrm{Fraction Shaded } = \frac{\answer{2}}{\answer{6}} \]

Which of the following circles represents an \textbf{equivalent} fraction?

%The * after hspace makes certain the space won't disappear if it happens to fall at the very beginning or very end of a line (the "boundary")
\hspace*{4em}  
\wordChoice{
    \choice{(a)}
    \choice[correct]{(b)}
    \choice{(c)}  
    }


\begin{center} 
\begin{tikzpicture}[thick, scale=0.5] 
% First Circle (a) 1/2 
\begin{scope}[xshift=0cm] 
\node at (0, 2.2) {\textbf{(a)}}; % Label (a) 
\fill[glaucous!60] (0,0) -- (90:1.3) arc (90:270:1.3) -- cycle; 
\draw (0,0) circle (1.3); 
\draw (0,1.3) -- (0,-1.3); 
\node at (0,-2.3) {$\frac{1}{2}$}; 
\end{scope} 

% Second Circle (b) 1/3 
\begin{scope}[xshift=5cm] % Increased shift for better spacing 
\node at (0, 2.2) {\textbf{(b)}}; % Label (b) 
% Changed the arc angles to 210 and 330 to shade the bottom slice
\fill[glaucous!60] (0,0) -- (210:1.3) arc (210:330:1.3) -- cycle; 
\draw (0,0) circle (1.3); 
\foreach \a in {90, 210, 330} \draw (0,0) -- (\a:1.3); 
\node at (0,-2.3) {$\frac{1}{3}$}; 
\end{scope} 

% Third Circle (c) 1/6 
\begin{scope}[xshift=10cm] 
\node at (0, 2.2) {\textbf{(c)}}; % Label (c) 
\fill[glaucous!60] (0,0) -- (30:1.3) arc (30:90:1.3) -- cycle; 
\draw (0,0) circle (1.3); 
\foreach \a in {30, 90, 150, 210, 270, 330} \draw (0,0) -- (\a:1.3); 
\node at (0,-2.3) {$\frac{1}{6}$}; 
\end{scope} 

\end{tikzpicture} 

\xmalt{Three circles labeled (a), (b), and (c) representing fractions. Circle (a) is divided into 2 equal parts with 1 shaded (one-half). Circle (b) is divided into 3 equal parts with 1 shaded (one-third). Circle (c) is divided into 6 equal parts with 1 shaded (one-sixth).} 

\end{center}

\end{problem}

\begin{definition}[\textbf{Types of Fractions}]

\begin{enumerate}
\item \textbf{Proper fractions} are fractions in which the numerator is smaller
than the denominator.  A proper fraction is always less than 1.

\begin{itemize}
        \item \textcolor{accessibleOrange}{\textbf{Example}}: \ $\frac{3}{4}$
\end{itemize}


\item \textbf{Improper fractions} are fractions in which the numerator is greater than
or equal to the denominator.  An improper fraction is always greater than or equal to 1.

\begin{itemize}
        \item \textcolor{accessibleOrange}{\textbf{Example}}: \ $\frac{7}{12}$
\end{itemize}


\item \textbf{Mixed numbers} consist of a whole number and a fraction.  A mixed number
is always greater than 1.

\begin{itemize}
        \item \textcolor{accessibleOrange}{\textbf{Example}}: \ $1 \frac{1}{4}$
\end{itemize}

\end{enumerate}

\end{definition}

\begin{problem}
Classify each of the following fractions as a \textbf{proper fraction}, an \textbf{improper fraction} or a \textbf{mixed number}.

\begin{enumerate}
    \item $\displaystyle 2\frac{1}{2}$
    \begin{multipleChoice}
        \choice{Proper Fraction}
        \choice{Improper Fraction}
        \choice[correct]{Mixed Number}
    \end{multipleChoice}

    \item $\displaystyle \frac{11}{12}$
    \begin{multipleChoice}
        \choice[correct]{Proper Fraction}
        \choice{Improper Fraction}
        \choice{Mixed Number}
    \end{multipleChoice}
    
    \item $\displaystyle \frac{15}{8}$
    \begin{multipleChoice}
        \choice{Proper Fraction}
        \choice[correct]{Improper Fraction}
        \choice{Mixed Number}
    \end{multipleChoice}

    \item $\displaystyle \frac{5}{5}$
    \begin{multipleChoice}
        \choice{Proper Fraction}
        \choice[correct]{Improper Fraction}
        \choice{Mixed Number}
    \end{multipleChoice}

\end{enumerate}
\end{problem}


\begin{problem}
When we write the mixed number $\textstyle 1 \frac{1}{4}$ what arithmetic symbol is not written, but
understood to be between the 1 and $\textstyle \frac{1}{4}?$

\begin{multipleChoice}
\choice{$\mathbf{\times}$}
\choice{$\mathbf{-}$}
\choice[correct]{$\mathbf{+}$}
\choice{$\mathbf{\div}$}
\end{multipleChoice}

\begin{feedback}[correct]
That's right! \ $1 \frac{1}{4} = 1 + \frac{1}{4}$
\end{feedback}

\end{problem}



\section*{Negative Fractions}

In the following problems, notice that the position of the negative sign changes but the value
of the fraction remains the same.

\begin{align*}
\frac{-4}{2} &= -2 \\[3ex]
\frac{4}{-2} &= -2 \\[3ex]
-\frac{4}{2} &= -2
\end{align*}

Therefore, a negative sign in a fraction can be written in the numerator, the denominator or in front
of the fraction.  However, when simplifying we typically place the negative sign in front of the fraction.


\end{document}
